\documentclass[11pt]{article}
\usepackage{amsmath, amssymb, amsthm}
\usepackage{geometry}
\usepackage{booktabs}
\usepackage{hyperref}
\usepackage{listings}
\usepackage{xcolor}
\usepackage{parskip}

\geometry{margin=1.25in}

\lstset{
  language=Python,
  basicstyle=\ttfamily\small,
  keywordstyle=\color{blue},
  commentstyle=\color{gray},
  stringstyle=\color{teal},
  frame=single,
  framexleftmargin=4pt,
  xleftmargin=8pt,
  breaklines=true,
}

\newtheorem{remark}{Remark}
\newcommand{\Nbasis}{N_{\mathrm{basis}}}
\newcommand{\Nosc}{N_{\mathrm{osc}}}
\newcommand{\ltol}{\texttt{ltol}}
\newcommand{\rtol}{\texttt{rtol}}
\newcommand{\ppl}{\texttt{ppl}}

\title{MPS Settings Strategy: Preliminary Heuristics}
\author{lappy}
\date{\today}

\begin{document}
\maketitle

\begin{abstract}
This document develops a principled basis for choosing \texttt{lappy} solver
settings---primarily \texttt{n\_basis}, \texttt{n\_fb}, \texttt{n\_fs},
\texttt{rtol}, \texttt{ltol}, and \texttt{ppl}---as functions of the desired
accuracy, the number of eigenvalues requested, and the geometry of the domain.
The analysis is preliminary; the benchmark sweeps are designed to refine the
constants and validate the functional forms.
\end{abstract}

\tableofcontents
\bigskip

%% ─────────────────────────────────────────────────────────────────────────────
\section{Background: the MPS tension function}

The Method of Particular Solutions (MPS) finds Laplacian eigenvalues as
approximate zeros of a \emph{tension function} $\sigma(\lambda)$.  Given a
basis of $N$ particular solutions $\{\varphi_1,\ldots,\varphi_N\}$---exact
solutions to $(\Delta+\lambda)u=0$ in $\mathbb{R}^2$---the tension at $\lambda$
is
\begin{equation}\label{eq:tension}
  \sigma(\lambda)
  =
  \min_{\|c\|=1}
  \frac{\bigl\|\sum_i c_i\varphi_i\bigr\|_{L^2(\partial\Omega)}}
       {\bigl\|\sum_i c_i\varphi_i\bigr\|_{L^2(\Omega)}}.
\end{equation}
$\sigma(\lambda)$ is small when there exists a linear combination that nearly
vanishes on $\partial\Omega$ while remaining non-trivial in the interior,
i.e.\ when $\lambda$ is near a true Dirichlet eigenvalue.  In the GSVD
formulation used here, $\sigma(\lambda)$ equals the smallest generalised
singular value of the pair (boundary evaluation matrix, interior evaluation
matrix).

The eigenvalue error is bounded by the tension:
\begin{equation}\label{eq:error_bound}
  \frac{|\hat\lambda - \lambda|}{\lambda}
  \;\lesssim\;
  C\,\sigma(\hat\lambda),
\end{equation}
where $C$ depends on the local spectral gap.  Achieving relative accuracy
$\ltol$ therefore requires $\sigma(\lambda)\lesssim\ltol$.

%% ─────────────────────────────────────────────────────────────────────────────
\section{Why convergence is exponential in \texorpdfstring{$N$}{N}}

The key theoretical fact underlying MPS is that the tension decays
\emph{exponentially} in $N$ for a well-chosen basis,
\begin{equation}\label{eq:exp_conv}
  \sigma(\lambda) \sim \exp(-\gamma N),
\end{equation}
where $\gamma>0$ is a geometry-dependent convergence rate.  This follows from
approximation theory for analytic functions.

\paragraph{Fourier--Bessel (FB) basis.}
At a corner with interior angle $\alpha$, the Dirichlet eigenfunctions have the
form
\begin{equation}\label{eq:fb_expansion}
  u(r,\theta) = \sum_{n\geq 1} a_n\, r^{n\pi/\alpha}\sin\!\left(\frac{n\pi\theta}{\alpha}\right)
\end{equation}
in local polar coordinates $(r,\theta)$ centred at the corner.  This is the
\emph{exact} singular form dictated by the geometry.  The FB basis at that
corner consists precisely of $J_{n\pi/\alpha}(\sqrt{\lambda}\,r)
\sin(n\pi\theta/\alpha)$, which matches this structure term by term.  The
residual after including $N$ terms is therefore the tail of an analytic series,
\begin{equation}
  \|u - u_N\| \sim C\exp(-\gamma_{\mathrm{corner}} N),
\end{equation}
where $\gamma_{\mathrm{corner}}$ is determined by the analyticity radius of $u$
away from that corner.  Since the eigenfunction is analytic everywhere except at
corners, and the FB basis handles the corner singularities exactly, the
remaining boundary residual falls off exponentially fast.

\paragraph{Fundamental Solution (FS) basis.}
Sources placed at distance $d$ outside the boundary generate functions that are
analytic on all of $\overline\Omega$.  For a smooth domain (or as a supplement
to FB functions on a polygon), $N$ such sources achieve
\begin{equation}
  \|u - u_N\| \sim C\exp(-\gamma_{\mathrm{fs}} N),
\end{equation}
where $\gamma_{\mathrm{fs}} \sim \pi d/L$ for sources distributed at distance
$d$ from a boundary of length $L$.  The convergence rate decreases as $d\to 0$,
so source placement matters.

\paragraph{Consequence for basis sizing.}
The tension target $\ltol$ requires
\begin{equation}\label{eq:N_accuracy}
  N \;\gtrsim\; \frac{|\log(\ltol)|}{\gamma}.
\end{equation}
This is a lower bound from the accuracy requirement alone.  As we will see, the
oscillation count often dominates in practice.

%% ─────────────────────────────────────────────────────────────────────────────
\section{The oscillation count argument}

The $n$-th eigenfunction oscillates at spatial frequency $\sqrt{\lambda_n}$.
Along a boundary segment of length $\ell$ it executes roughly
$\ell\sqrt{\lambda_n}/(2\pi)$ oscillations.  A basis of $N$ functions centred
at a single point can accurately represent at most $O(N)$ oscillation modes.
To capture the first $k$ eigenvalues across all boundary segments we need the
basis to span the oscillations at scale $\sqrt{\lambda_k}$.  Summing over all
corners or source locations,
\begin{equation}\label{eq:Nosc}
  \Nosc \sim C\,\frac{P\sqrt{\lambda_k}}{2\pi},
\end{equation}
where $P$ is the perimeter.  Substituting Weyl's law
$\lambda_k \sim 4\pi k/A$ gives $\sqrt{\lambda_k} \sim 2\sqrt{\pi k/A}$, and
hence
\begin{equation}\label{eq:Nosc_weyl}
  \Nosc \sim C\,\frac{P}{\sqrt{A}}\,\sqrt{k}.
\end{equation}
The factor $P/\sqrt{A}$ is the \emph{isoperimetric ratio}.  It equals
$2\sqrt\pi\approx 3.54$ for a circle and grows for elongated or geometrically
complex domains.  Representative values: a square has $P/\sqrt{A}=4$; an
$L$-shape has $P/\sqrt{A}\approx 5$--$6$; the GWW isospectral domains have
$P/\sqrt{A}\approx 7$--$9$.

%% ─────────────────────────────────────────────────────────────────────────────
\section{The combined \texorpdfstring{$\Nbasis$}{N\_basis} heuristic}

The two constraints---accuracy \eqref{eq:N_accuracy} and oscillation count
\eqref{eq:Nosc_weyl}---combine as
\begin{equation}\label{eq:heuristic}
  \Nbasis
  \;\sim\;
  C\,\frac{P}{\sqrt{A}}\,\sqrt{k}\;F_{\mathrm{geom}}\;f(\ltol),
\end{equation}
where we now describe each factor.

\subsection{Accuracy factor \texorpdfstring{$f(\ltol)$}{f(ltol)}}

Once $N$ exceeds the oscillation count, the tension drops rapidly.  Because
convergence is exponential, each additional basis function buys a fixed
multiplicative improvement.  A reasonable approximation is
\[
  f(\ltol) = \max\!\left(1,\;\alpha\,|\log(\ltol)|^\beta\right),
  \quad \beta\in[0.5,\,1.0].
\]
The benchmark data shows that going from $\ltol=10^{-10}$ to $10^{-16}$ requires
only 30--50\% more basis functions, consistent with $\beta\approx 0.5$.  We
conservatively use $\beta=1$ (linear in $|\log(\ltol)|$) until more data are
available.

\subsection{Geometric difficulty factor \texorpdfstring{$F_{\mathrm{geom}}$}{F\_geom}}

\paragraph{Re-entrant corners ($\alpha>\pi$).}
The singularity exponent $\pi/\alpha<1$ means the eigenfunction is less smooth
at the corner; more FB terms are needed to resolve it.  A rough per-corner
factor is $(\alpha/\pi)^{1/2}$ for $\alpha>\pi$, giving
\begin{equation}
  F_{\mathrm{corners}}
  = \prod_{\substack{i\,:\,\text{corner }i\\\alpha_i>\pi}}
    \left(\frac{\alpha_i}{\pi}\right)^{1/2}.
\end{equation}
For example, the $L$-shape has one $270^\circ$ corner ($\alpha=3\pi/2$), giving
$F_{\mathrm{corners}}=\sqrt{3/2}\approx 1.22$.  This is consistent with
needing $N\approx 120$ for the $L$-shape versus $N\approx 60$--$80$ for
rectangles and triangles.

\paragraph{Irrational corner angles.}
If $\pi/\alpha_i$ is irrational, the required FB orders are non-integer-order
Bessel functions with slower convergence than integer-order ones.  The disk
sector at $\theta=\pi\sqrt{2}/2$ is the clearest example in the benchmark data;
it struggles even at $\Nbasis=160$.  For irrational angles the convergence
rate $\gamma$ is degraded and a larger empirical multiplier is needed.

\paragraph{Smooth curved boundaries.}
Domains with curved sides but no corners (disk, ellipse, stadium) require no FB
functions.  The FS convergence rate depends on the source placement distance
$d$ relative to the inradius $r_{\mathrm{in}}$:
$\gamma_{\mathrm{fs}} \sim \pi d/r_{\mathrm{in}} / P_{\mathrm{norm}}$.
A source distance of $d\approx 0.15\,r_{\mathrm{in}}$ is the current default;
the \texttt{fs\_placement} sweep will quantify the optimal value.

\subsection{Summary and comparison with data}

Putting it together, the working preliminary heuristic is
\begin{equation}\label{eq:heuristic_final}
  \Nbasis
  \;\approx\;
  \mathrm{round}\!\left(
    C\,\frac{P}{\sqrt{A}}\,\sqrt{k}\;F_{\mathrm{geom}}\;|\log_{10}(\ltol)|
  \right),
\end{equation}
with $C\approx 2$--$4$ to be fit from the benchmark data.  Table~\ref{tab:predictions}
compares predictions to preliminary observations.

\begin{table}[h]
\centering
\begin{tabular}{lcccc}
\toprule
Domain & $P/\sqrt{A}$ & $F_{\mathrm{geom}}$ & Predicted $N$ ($k=10$, $\ltol=10^{-14}$) & Observed $N$ \\
\midrule
rect ($2\times 1$)  & 4.2 & 1.0  & 55--70   & $\approx 60$  \\
iso\_right\_tri     & 4.6 & 1.0  & 60--75   & $\approx 70$  \\
eq\_tri             & 4.6 & 1.0  & 60--75   & $\approx 60$  \\
L\_shape            & 5.5 & 1.22 & 85--110  & $\approx 120$ \\
GWW1                & $\approx 8$ & $\approx 1.5$ & 150--200 & $>160$ (tbc) \\
\bottomrule
\end{tabular}
\caption{Predicted versus observed basis sizes required for $k=10$ eigenvalues
at relative accuracy $\ltol=10^{-14}$.  Predictions use
\eqref{eq:heuristic_final} with $C=3$.}
\label{tab:predictions}
\end{table}

%% ─────────────────────────────────────────────────────────────────────────────
\section{Role of \texorpdfstring{$k$}{k} (number of requested eigenvalues)}

The number of requested eigenvalues $k$ enters through the $\sqrt{k}$ scaling
of the oscillation count in \eqref{eq:Nosc_weyl}.  Key implications:

\paragraph{Sublinear cost growth.}
Doubling $k$ requires only $\sqrt{2}\approx 41\%$ more basis functions, not
double.  This is one of the attractive properties of MPS relative to finite
element methods, which typically scale as $O(k)$ or worse in practice.

\paragraph{Sizing for the largest eigenvalue.}
The basis must be sized for $\lambda_k$, the \emph{largest} eigenvalue requested.
Using a basis sized for $k=10$ when $k=25$ eigenvalues are requested will
reliably find the first few but miss or inaccurately compute the upper range.
This is open issue~2.7 (adaptive basis sizing): currently \texttt{evp.solve(k)}
uses whatever basis was constructed at initialization regardless of $k$.

\paragraph{Practical lower bound.}
\begin{equation}
  \Nbasis \;\geq\; N_{\min}(k)
  = \mathrm{round}\!\left(C\,\frac{P}{\sqrt{A}}\,\sqrt{k}\;F_{\mathrm{geom}}\right),
\end{equation}
with $C\approx 4$--$6$ as a conservative starting point.  The GSVD cost scales
as $O(N^2 n_{\mathrm{pts}})$, so there is a penalty for over-sizing, but it is
not catastrophic for moderate $N$.

\paragraph{Interaction with \ppl.}
The \ppl\ parameter (points per level in the bracket search) also interacts with
$k$.  The solver searches for $k$ tension minima on $[\lambda_{\min},\lambda_{\max}]$.
With too few points per level (\texttt{ppl}$=5$), eigenvalues can be missed
entirely---this is the primary cause of solver failures in the benchmark data at
small $\Nbasis$.  The recommended minimum is \texttt{ppl}$=10$; for difficult
domains or large $k$, \texttt{ppl}$=15$--$20$ is safer.

%% ─────────────────────────────────────────────────────────────────────────────
\section{The FB/FS split}

Given a total budget of $\Nbasis$ functions, how should they be divided between
Fourier-Bessel (\texttt{n\_fb}) and Fundamental Solutions (\texttt{n\_fs})?

\paragraph{Pure polygon, no curved sides.}
FB only (\texttt{n\_fs}$=0$).  The FB basis exactly matches the singularity
structure; FS functions add smooth global corrections that are redundant once the
FB orders are high enough.

\paragraph{Pure smooth domain (no corners).}
FS only (\texttt{n\_fb}$=0$).  There is no singularity structure for the FB
basis to capture.

\paragraph{Mixed domain (polygon with curved sides).}
Use both.  FB functions handle corner singularities; FS functions supplement the
smooth regions.

\paragraph{FB order allocation across corners.}
With $n_{\mathrm{fb}}$ functions distributed across $n_c$ corners, the
\texttt{fb\_strategy} parameter controls allocation:
\begin{itemize}
  \item \texttt{singular\_angle\_weighted} (default): allocates more orders to
    corners with larger (more re-entrant) angles.  This is theoretically
    motivated since more singular corners require more FB terms.
  \item \texttt{uniform}: equal orders per corner.  Suboptimal when corner
    angles vary widely.
  \item \texttt{singular\_only}: all orders at non-regular corners (those where
    $\pi/\alpha$ is not an integer), zero at regular corners.  Aggressive but
    correct in theory for polygons with a mix of regular and singular corners.
\end{itemize}
The \texttt{fb\_corner\_sweep} will quantify the accuracy difference between
strategies.  The preliminary expectation is that for domains with one dominant
singular corner (such as the $L$-shape), \texttt{singular\_only} and
\texttt{singular\_angle\_weighted} will outperform \texttt{uniform}.

%% ─────────────────────────────────────────────────────────────────────────────
\section{Collocation points}

The boundary and interior collocation points affect accuracy and stability, but
are less sensitive than basis size.

\paragraph{Boundary points (\texttt{bdry\_pts\_factor}).}
The linear system is overdetermined: we use
$\texttt{bdry\_pts\_factor}\times\Nbasis$ boundary points.  A factor of $2.0$
(the default) is standard for MPS and gives good conditioning.  Going below
$1.5$ risks an underdetermined system; going above $3.0$ adds cost without
meaningful benefit.  Points should be distributed proportionally to segment
length.  Gauss--Legendre quadrature points are preferable to uniform spacing
because they avoid endpoint clustering artifacts.

\paragraph{Interior points (\texttt{int\_pts\_factor}).}
Used for basis normalisation ($L^2$ norm over $\Omega$).  A factor of $1.0$ is
typically sufficient; the normalisation is not sensitive to this choice as long
as the points are reasonably distributed.

%% ─────────────────────────────────────────────────────────────────────────────
\section{Regularisation: \texttt{rtol} and \texttt{ltol}}

The GSVD-based MPS involves two tolerance parameters with different roles.

\paragraph{\rtol\ (regularisation tolerance).}
Singular values of the basis matrix below $\rtol\times\sigma_{\max}$ are
discarded.  This controls ill-conditioning from near-linearly-dependent basis
functions.
\begin{itemize}
  \item Too small: numerical noise dominates and the tension function becomes
    noisy, causing missed or spurious eigenvalues.
  \item Too large: legitimate but small singular values (corresponding to weakly
    excited modes) are discarded, causing missed eigenvalues.
\end{itemize}
Preliminary results from the \texttt{regularization\_sweep} suggest the sweet
spot is $\rtol\in[10^{-12},10^{-14}]$ for well-conditioned domains.  The GWW
domain degrades significantly below $\rtol=10^{-12}$, suggesting the basis
becomes ill-conditioned at large $N$ for complex geometries.

A geometry-motivated heuristic is
\begin{equation}
  \rtol \;\sim\; C_r\exp(-\gamma\,\Nbasis),
\end{equation}
i.e., set \rtol\ to the expected tension floor for the given basis size and
geometry.  In practice $\rtol=10^{-14}$ is a safe default; tighten to
$10^{-16}$ only if the basis is verified to be well-conditioned.

\paragraph{\ltol\ (eigenvalue acceptance threshold).}
A tension minimum is accepted as an eigenvalue only if $\sigma(\hat\lambda)<\ltol$.
This should be set relative to the accuracy the basis is expected to achieve:
\begin{equation}
  \ltol \;\sim\; 10\times\exp(-\gamma\,\Nbasis) \;\sim\; 10\times\sigma_{\mathrm{floor}}(\Nbasis).
\end{equation}
Setting $\ltol$ too tight causes real eigenvalues to be rejected; too loose
accepts noise peaks as eigenvalues.  The current default of $\ltol=10^{-15}$ is
aggressive; for small bases ($N<50$) or complex domains,
$\ltol=10^{-10}$--$10^{-12}$ is more reliable.

\paragraph{The \rtol/\ltol\ relationship.}
Both tolerances should be set consistently.  A rough guideline is
\begin{equation}
  \ltol \;\approx\; 100\times\rtol,
\end{equation}
which ensures that basis regularisation does not create artificial tension minima
above the eigenvalue acceptance threshold.

%% ─────────────────────────────────────────────────────────────────────────────
\section{Practical recommendations (preliminary)}

Until the benchmark sweeps are complete, the following settings are conservative
starting points.

\paragraph{Step 1: estimate $\Nbasis$.}

\begin{lstlisting}
import numpy as np

def estimate_n_basis(domain, k, ltol=1e-14):
    P = domain.perim
    A = domain.area
    iso_ratio = P / np.sqrt(A)
    # geometric difficulty from re-entrant corners
    angles = domain.int_angles
    F_geom = np.prod([np.sqrt(a / np.pi) for a in angles if a > np.pi]) or 1.0
    C = 5.0   # conservative constant; refine from benchmarks
    return round(C * iso_ratio * np.sqrt(k) * F_geom * abs(np.log10(ltol)))
\end{lstlisting}

\paragraph{Step 2: set the FB/FS split.}
\begin{itemize}
  \item Polygon with all straight sides: \texttt{n\_fs}$=0$, \texttt{n\_fb}$=\Nbasis$.
  \item Smooth domain: \texttt{n\_fb}$=0$, \texttt{n\_fs}$=\Nbasis$.
  \item Mixed: \texttt{n\_fb}$=\mathrm{round}(0.5\,\Nbasis)$,
    \texttt{n\_fs}$=\Nbasis-\texttt{n\_fb}$.
\end{itemize}

\paragraph{Step 3: set tolerances.}

\begin{table}[h]
\centering
\begin{tabular}{lccc}
\toprule
Use case & \rtol & \ltol & \ppl \\
\midrule
Quick / exploratory         & $10^{-10}$ & $10^{-8}$  & 10 \\
Standard                    & $10^{-12}$ & $10^{-10}$ & 10 \\
High accuracy               & $10^{-14}$ & $10^{-12}$ & 15 \\
Near-degenerate eigenvalues & $10^{-14}$ & $10^{-12}$ & 20 \\
\bottomrule
\end{tabular}
\caption{Recommended tolerance settings by use case.}
\end{table}

\paragraph{Step 4: check the result.}
A result should be treated as reliable if:
\begin{itemize}
  \item \texttt{len(eigs) == k} (all requested eigenvalues found),
  \item \texttt{max(tensions) < ltol} (all eigenvalues accepted with margin),
  \item \texttt{median(rel\_errors) < ltol} (when reference eigenvalues are
    available).
\end{itemize}
If \texttt{len(eigs) < k}, increase $\Nbasis$ by 25--50\% and/or increase
\ppl.

%% ─────────────────────────────────────────────────────────────────────────────
\section{Open questions for the benchmark data}

The following will be addressed as more sweep results become available.

\begin{enumerate}
  \item \textbf{The constant $C$.}  The current range $C\approx 2$--$6$ is too
    wide.  The \texttt{basis\_size} sweep will pin this down per domain class.

  \item \textbf{The accuracy exponent $\beta$.}  Is $f(\ltol)\sim|\log(\ltol)|$
    or $\sim|\log(\ltol)|^{1/2}$?  The \texttt{regularization\_sweep} data will
    answer this.

  \item \textbf{Optimal FB/FS ratio for mixed domains.}  The
    \texttt{basis\_size} sweep covers \texttt{fb\_fraction}$\in\{0,\,0.25,\,0.5,\,0.75,\,1.0\}$;
    we expect a clear optimum around $0.5$--$0.75$ for polygons.

  \item \textbf{FS source distance scaling.}  The \texttt{fs\_placement} sweep
    will determine how accuracy depends on \texttt{fs\_d\_scale}.  The current
    theoretical prediction is $\gamma_{\mathrm{fs}}\sim\pi d/r_{\mathrm{in}}$;
    the data will test this.

  \item \textbf{The irrational angle case.}  The disk sector at
    $\theta=\pi\sqrt{2}/2$ is a red flag.  We need to understand whether this
    is a basis sizing issue, an FB order allocation issue, or a fundamental
    limitation of the FB convergence rate.

  \item \textbf{The $\sqrt{k}$ scaling exponent.}  The $\sqrt{k}$ prediction
    from Weyl's law should be directly testable from the \texttt{n\_eigs} sweep.
\end{enumerate}

\end{document}
